\documentclass[11pt]{article}
\usepackage[margin=1in]{geometry}
\usepackage{amsmath,amssymb,amsthm}
\usepackage{booktabs}
\usepackage{hyperref}

\newtheorem{theorem}{Theorem}
\newtheorem{lemma}[theorem]{Lemma}
\newtheorem{proposition}[theorem]{Proposition}
\newtheorem{corollary}[theorem]{Corollary}
\theoremstyle{definition}
\newtheorem{definition}[theorem]{Definition}
\theoremstyle{remark}
\newtheorem{remark}[theorem]{Remark}

\newcommand{\F}{\mathcal F}
\newcommand{\ip}[2]{\langle #1,#2\rangle}
\newcommand{\NC}{\mathrm{NC}_2}

\title{Disproof of Conjecture 00000007739:\\
the fourth moment of a free semicircular word of length one is $2$, not $9/4$}
\author{jilint777}
\date{2026-10-04}

\begin{document}
\sloppy
\maketitle

\begin{abstract}
Conjecture 00000007739 asserts that for free words $W$ of length $d$ in a free semicircular
family, the ratio $M_p(d)=\|W\|_p/\|W\|_2$ satisfies
$M_4(d)^2=2-d^{-1}(1-2^{-d})$. This is false already for $d=1$. A word of length one is a single
standard semicircular element $s$, with $\tau(s^2)=1$ and $\tau(s^4)=2$, so $M_4(1)^4=2$ and
$M_4(1)^2=\sqrt2$, whereas the formula gives $3/2$ and $(3/2)^2=9/4\ne2$. The same value
$M_4(1)^4=2$ holds for \emph{every} nonzero element of the first semicircular chaos (real or
complex coefficients), so reading $M_p(d)$ as an optimal (or extremal) constant changes nothing.
More generally, for $W=s_1\cdots s_d$ we have $M_4(d)^4=d+1$, and the formula fails for every
$d\ge1$ and under each of the normalisations $M_4$, $M_4^2$, $M_4^4$. The decisive objects
(Voiculescu's semicircular elements $\ell(e_i)+\ell(e_i)^*$ on the full Fock space, the vacuum
state, adjoints and the norms) are constructed from scratch in Lean~4 (core library only), and the
negation of the clause is proved there. An independent Python script recomputes all moments
exactly by two methods.
\end{abstract}

\section{The claim}

The conjecture (English version): \emph{``The free Khintchine inequalities compare $L^p$ and $L^2$
norms of noncommutative martingale polynomials: for free words $W$ of length $d$ generated by free
semicircular families, $M_p(d)=\|W\|_p/\|W\|_2$. Conjecture: $M_4(d)^2=2-d^{-1}(1-(1/2)^d)$, the
exact free refinement of the fourth-moment theorem; and the limit ratio spectrum of
$M_\infty(d)/M_p(d)$ as $d$ grows characterizes the operator-norm concentration of products of free
group generators.''} The Chinese version calls $M_p(d)$ ``the optimal constant of the ratio''.

The statement is a conjunction of
\begin{itemize}
\item[(C1)] $M_4(d)^2=f(d)$ for every $d\ge1$, where
\[
f(d):=2-\frac{1-2^{-d}}{d}=\frac{2d\cdot2^d-2^d+1}{d\cdot2^d},
\]
so that $f(1)=\tfrac32$, $f(2)=\tfrac{13}8$, $f(3)=\tfrac{41}{24}$, $f(4)=\tfrac{113}{64}$ and
$f(5)=\tfrac{289}{160}$;
\item[(C2)] an informal statement on $M_\infty(d)/M_p(d)$ and free group generators.
\end{itemize}
We refute (C1); hence the conjunction is false whatever (C2) means. Note $\tfrac32\le f(d)<2$ for
all $d\ge1$ (since $0<(1-2^{-d})/d\le\tfrac12$).

\section{Set-up}

Let $(\mathcal A,\tau)$ be a tracial $W^*$-probability space. For $x\in\mathcal A$ and $p<\infty$,
$\|x\|_p=\tau(|x|^p)^{1/p}$ with $|x|=(x^*x)^{1/2}$. In particular
\[
\|x\|_2^2=\tau(x^*x),\qquad \|x\|_4^4=\tau\big((x^*x)^2\big).
\]
The ratio $\|x\|_4/\|x\|_2$ is invariant under $x\mapsto cx$ ($c\neq0$).

\paragraph{Full Fock space (Voiculescu).} Let $H$ be a Hilbert space with orthonormal basis
$(e_i)_{i\ge0}$ and $\F(H)=\mathbb C\Omega\oplus\bigoplus_{n\ge1}H^{\otimes n}$, with orthonormal
basis $e_{i_1}\otimes\cdots\otimes e_{i_n}$ (the empty word is the vacuum $\Omega$). The creation
operator $\ell(e_i)\xi=e_i\otimes\xi$ has adjoint $\ell(e_i)^*$, given by $\ell(e_i)^*\Omega=0$ and
$\ell(e_i)^*(e_j\otimes\xi)=\delta_{ij}\xi$. Put
\[
X_i=\ell(e_i)+\ell(e_i)^*,\qquad \tau(T)=\ip{\Omega}{T\Omega}.
\]
By Voiculescu's theorem, $X_0,X_1,\dots$ are free standard semicircular elements with respect to
$\tau$, and $\tau$ is a faithful trace on the von Neumann algebra they generate. Since any two free
semicircular families of the same size have the same joint $*$-distribution, every quantity
below is the same in every realisation. For a self-adjoint $y$ in this algebra,
$\tau(y^2)=\|y\Omega\|^2$. Hence, for every $x$,
\begin{equation}\label{eq:norms}
\|x\|_2^2=\|x\Omega\|^2,\qquad \|x\|_4^4=\|x^*x\,\Omega\|^2 .
\end{equation}

\paragraph{Moment formula.} For the free semicircular family,
$\tau(X_{i_1}\cdots X_{i_m})$ is the number of non-crossing pair partitions $\pi\in\NC(m)$ such
that $i$ is constant on each block of $\pi$ (Speicher's free Wick formula; it also follows directly
by expanding the product of $\ell$'s and $\ell^*$'s applied to $\Omega$).

\section{The decisive case \texorpdfstring{$d=1$}{d=1}}

\begin{theorem}\label{thm:main}
Let $s$ be a standard semicircular element. Then $\|s\|_2^2=1$ and $\|s\|_4^4=2$, so
$M_4(1)^4=2$. Consequently $M_4(1)^2=\sqrt2\ne\tfrac32=f(1)$, and clause (C1) is false.
\end{theorem}

\begin{proof}
Realise $s=X_0$. Then $X_0\Omega=e_0$ and
$X_0^2\Omega=X_0e_0=e_0\otimes e_0+\Omega$. By \eqref{eq:norms},
$\|s\|_2^2=\|e_0\|^2=1$ and $\|s\|_4^4=\|e_0\otimes e_0+\Omega\|^2=1+1=2$.
(Equivalently, there are two non-crossing pairings of four points.) If $M_4(1)^2=f(1)=3/2$, then
squaring gives $2=M_4(1)^4=9/4$, which is false.
\end{proof}

A word of length one in a free semicircular family is a single $s_i$, so for $d=1$ there is no
ambiguity about which word is meant.

\section{Optimal-constant reading: the whole first chaos}

The Chinese text suggests that $M_p(d)$ is the \emph{optimal constant}, i.e.\ the supremum of
$\|W\|_p/\|W\|_2$ over a class of degree-$d$ elements. For $d=1$ the natural class is the first
chaos $\{\sum_i c_iX_i\}$ (finitely many nonzero $c_i$, or $c\in\ell^2$ by continuity), with real or
complex coefficients.

\begin{proposition}\label{prop:chaos1}
For every nonzero $S=\sum_i c_iX_i$ with $c_i\in\mathbb C$: $\|S\|_4^4=2\|S\|_2^4$. Hence the ratio
$\|S\|_4/\|S\|_2=2^{1/4}$ is constant on the first chaos. Its supremum, infimum and every value it
takes give $M_4(1)^4=2$.
\end{proposition}

\begin{proof}
Let $\beta=\sum_i|c_i|^2$. Then $S\Omega=\sum_jc_je_j=:\xi$, so $\|S\|_2^2=\beta$. Next
$S^*=\sum_i\bar c_iX_i$, and $\ell(e_i)^*\xi=c_i\Omega$, so
\[
S^*S\Omega=\sum_i\bar c_i\big(e_i\otimes\xi+c_i\Omega\big)
=\sum_{i,j}\bar c_ic_j\,e_i\otimes e_j+\beta\,\Omega .
\]
The two summands are orthogonal, so
$\|S^*S\Omega\|^2=\sum_{i,j}|c_i|^2|c_j|^2+\beta^2=2\beta^2$, and \eqref{eq:norms} finishes the
proof.
\end{proof}

For real coefficients this is the stability of the semicircle law: $\sum c_iX_i$ is semicircular
of variance $\sum c_i^2$. For complex coefficients the circular element
$(X_0+iX_1)/\sqrt2$ and the elliptic elements $aX_0+ibX_1$ are covered. Thus under the
optimal-constant reading, real or complex, $M_4(1)^2=\sqrt2\ne3/2$ again.

\section{General \texorpdfstring{$d$}{d}}

\begin{proposition}\label{prop:fuss}
Let $W_d=s_1s_2\cdots s_d$ with $s_1,\dots,s_d$ free standard semicirculars. Then
$\|W_d\|_2^2=1$ and $\|W_d\|_4^4=d+1$.
\end{proposition}

\begin{proof}
$\tau(W_d^*W_d)=\tau(s_d\cdots s_1s_1\cdots s_d)$. Every colour occurs twice, and the only
admissible pairing pairs the two occurrences of each colour. This pairing is non-crossing (the
pairs are nested), so the moment is $1$.

For $\tau((W_d^*W_d)^2)$, write the word as $A\,B\,C\,D$ with $A=C=(d,\dots,1)$ and
$B=D=(1,\dots,d)$. Colour $k$ occurs at positions $a_k=d+1-k$, $b_k=d+k$, $c_k=3d+1-k$ and
$d_k=3d+k$. A monochromatic pairing pairs, for each $k$, these four points as
$(a_kb_k)(c_kd_k)$ (``inner''), $(a_kd_k)(b_kc_k)$ (``outer''), or $(a_kc_k)(b_kd_k)$. The last
is crossing. For $k<l$:
\begin{itemize}
\item if $k$ is outer and $l$ is inner, then $a_l<a_k<b_l<d_k$, so $(a_lb_l)$ crosses $(a_kd_k)$;
\item in the three other cases the pairs are nested or disjoint. If both are inner, $(a_kb_k)$ lies
inside $(a_lb_l)$ and $(c_kd_k)$ inside $(c_ld_l)$. If $k$ is inner and $l$ is outer,
$(a_kb_k)$ and $(c_kd_k)$ lie inside $(a_ld_l)$ and outside $(b_lc_l)$. If both are outer,
$(a_kd_k)\supset(b_lc_l)$ and $(b_kc_k)\supset(b_lc_l)$, all inside $(a_ld_l)$.
\end{itemize}
So the non-crossing monochromatic pairings correspond to the sets of outer colours that are
up-sets $\{m,\dots,d\}$ of $\{1,\dots,d\}$, $m\in\{1,\dots,d+1\}$. There are $d+1$ of them.
(This is the Fuss--Catalan number $\frac1{2d+1}\binom{2d+2}{2}$, the second moment of the free
multiplicative convolution of $d$ free Poisson laws.)
\end{proof}

\begin{corollary}\label{cor:alld}
For every $d\ge1$ the word $W_d$ violates (C1), and also the variants $M_4(d)=f(d)$ and
$M_4(d)^4=f(d)$.
\end{corollary}

\begin{proof}
$M_4(d)^4=d+1$ is an integer. Each variant, raised to the appropriate power, would give
$d+1=f(d)^k$ with $k\in\{1,2,4\}$. But the numerator $2d\cdot2^d-2^d+1$ of $f(d)$ is odd, so the
$2$-adic valuation is $v_2(f(d))=-d-v_2(d)<0$, and $f(d)^k$ is never an integer.
\end{proof}

Directly: $M_4(d)^2=\sqrt{d+1}\ge\sqrt2$, while $f(d)<2$. Equality would require
$d+1=f(d)^2\in[9/4,4)$, i.e.\ $d=2$, but $f(2)^2=169/64\ne3$.

\section{Readings covered}

\begin{itemize}
\item \textbf{Normalisation.} At $d=1$: $M_4(1)=2^{1/4}$, $M_4(1)^2=\sqrt2$ and $M_4(1)^4=2$.
The formula value $3/2$ matches none of them: $(3/2)^4=81/16\ne2$, $(3/2)^2=9/4\ne2$ and
$3/2\ne2$. For general $d$ see Corollary~\ref{cor:alld}.
\item \textbf{Which word.} For $d=1$ every word is a single semicircular. For $d\ge2$ we use the
word with distinct letters, $s_1\cdots s_d$. Words with repeated letters give other values (e.g.\
$\|s^2\|_4^4/\|s^2\|_2^4=14/4$), so $M_p(d)$ is then only meaningful as an extremal constant.
\item \textbf{Optimal constant (sup), or inf, over the $d$-th chaos; real or complex
coefficients.} For $d=1$ the ratio is constantly $2$ (Proposition~\ref{prop:chaos1}). For $d\ge2$:
the $d$-th Wigner chaos is $\{x:\ x\Omega\in H^{\otimes d}\}$. For such $x$ and
$\xi\in H^{\otimes d}$, the free Wick formula
$x=\sum_k\ell(\cdot)^k\ell(\cdot)^{*\,d-k}$ (Biane--Speicher) shows that the
$H^{\otimes 2d}$-component of $x^*\xi$ is $(x^*\Omega)\otimes\xi$ and its $\Omega$-component is
$\ip{x\Omega}{\xi}\Omega$. Taking $\xi=x\Omega$ and using traciality gives
$\|x\|_4^4=\|x^*x\Omega\|^2\ge2\|x\|_2^4$ (for self-adjoint $x$ this is the bound of
Kemp--Nourdin--Peccati--Speicher). So the ratio $\|x\|_4^4/\|x\|_2^4$ is $\ge2$, its supremum is
$\ge d+1$ (Proposition~\ref{prop:fuss}), and its infimum is exactly $2$. Indeed, for $n$ free copies
$w_1,\dots,w_n$ of $s_1\cdots s_d$, the element $\sum_jw_j$ has ratio $2+(d-1)/n$
(\texttt{verify.py} checks this). Since $9/4\le f(d)^2<4$ and $f(2)^2=169/64<3$, neither the
supremum nor the infimum equals $f(d)^2$ for any $d\ge1$. Since $f(d)<2$, no element has
ratio $f(d)$ (variant $M_4^4=f$). In the variant $M_4=f$, the infimum $2$ differs from
$f(d)^4\ge81/16$. We did not determine the exact supremum for $2\le d\le14$ in that variant.
This is not needed, because $d=1$ already refutes the clause in every variant.
\item \textbf{Haar unitaries / free group generators} (suggested by (C2)). A word in free Haar
unitaries $u_i=\lambda(g_i)$ is unitary, so $|W|=1$ and $M_p(d)=1$ for all $p$; $1\ne f(d)$,
$f(d)^2$, $f(d)^4$. For linear combinations $x=\sum_sa_s\lambda(s)$ over generators and their
inverses, the identity $s^{-1}tu^{-1}v=e$ holds iff ($s=t$, $u=v$) or ($t=u$, $s=v$). Hence
\[
\frac{\|x\|_4^4}{\|x\|_2^4}=2-\frac{\sum_s|a_s|^4}{\big(\sum_s|a_s|^2\big)^2}\in[1,2).
\]
So $M_4^2<\sqrt2<3/2$ and $M_4<3/2$. In the variant $M_4^4=f(1)=3/2$, this value \emph{is} attained,
by $u+u^*$ (the arcsine law: $\tau((u+u^*)^4)/\tau((u+u^*)^2)^2=6/4$). But it is not the optimal
constant, which is $\sup=2$ (or $\inf=1$). Moreover the conjecture explicitly concerns
\emph{semicircular} families, where the value is $2$.
\end{itemize}

\section{Machine verification}

\paragraph{Lean 4 (\texttt{lean4/Main.lean}, core only, v4.19.0).} Everything is defined from
scratch:
\begin{itemize}
\item basis words \texttt{Word := List Nat} of the full Fock space, and finitely supported
vectors \texttt{Vec := List (Word $\times$ Int)} (formal finite combinations; \texttt{coeff},
\texttt{inner}). All vectors that occur have integer coordinates.
\item \texttt{create i} $=\ell(e_i)$, \texttt{annih i} $=\ell(e_i)^*$, \texttt{X i}
$=\ell(e_i)+\ell(e_i)^*$, the vacuum \texttt{$\Omega$}, and the state
\texttt{tau ws} $=\ip{\Omega}{X_{i_1}\cdots X_{i_k}\Omega}$.
\item adjointness: \texttt{inner\_create}, \texttt{inner\_annih} (so $\ell(e_i)^*$ is the adjoint
of $\ell(e_i)$ on finitely supported vectors), \texttt{X\_selfadjoint}, and
\texttt{applyW\_adjoint}: the adjoint of $X_{i_1}\cdots X_{i_k}$ is $X_{i_k}\cdots X_{i_1}$.
\item \texttt{norm2sq} $=\tau(W^*W)$ and \texttt{norm4pow4} $=\tau((W^*W)^2)$, with
\texttt{norm2sq\_eq}, \texttt{norm4pow4\_eq} proving \eqref{eq:norms}.
\end{itemize}
Results:
\begin{itemize}
\item \texttt{semicircle\_moments}: $\tau(X_0^k)$ for $k\le10$ are $1,0,1,0,2,0,5,0,14,0,42$.
\item \texttt{mixed\_moments}: sample mixed moments of $X_0,X_1$.
\item \texttt{word1}: $\|X_0\|_2^2=1$ and $\|X_0\|_4^4=2$.
\item \texttt{words\_234}: $\|X_0\cdots X_{d-1}\|_4^4=d+1$ and $\|\cdot\|_2^2=1$ for $d=2,3,4$.
\item \texttt{first\_chaos\_three}: for all integers $a,b,c$, $S=aX_0+bX_1+cX_2$ satisfies
$\tau(S^2)=a^2+b^2+c^2$ and $\tau(S^4)=2(a^2+b^2+c^2)^2$. This is a polynomial identity, so it
holds for real coefficients too.
\item \texttt{first\_chaos\_never\_formula}: no nonzero such $S$ has ratio $9/4$, $3/2$ or
$81/16$.
\item \texttt{Clause d} encodes $M_4(d)^2=f(d)$ as $N(d)^2\|W_d\|_2^4=\|W_d\|_4^4D(d)^2$, where
$f=N/D$. Since $f(d)>0$, $f(d)=\sqrt{y}$ is equivalent to $f(d)^2=y$, so no real numbers are needed.
\item \texttt{clause\_false\_d1}: $\neg$\,\texttt{Clause 1}.
\item \texttt{clause\_false\_small}: \texttt{Clause} and the variants \texttt{ClauseM4},
\texttt{ClauseM4pow4} fail for $d=1,\dots,4$.
\item \texttt{conjecture\_00000007739\_false (C2) :} $\neg(\texttt{Conj}\wedge C_2)$, where
\texttt{Conj} $:=\forall d\ge1$, \texttt{Clause d}.
\end{itemize}
There is no \texttt{sorry} and no \texttt{native\_decide}. The axioms used are at most
\texttt{propext} and \texttt{Quot.sound}.

\emph{Limitations.} Lean works on the finitely supported vectors of the Fock space; all vectors
reached from $\Omega$ are finitely supported. The following are cited or proved in this report
only:
\begin{itemize}
\item the passage to the von Neumann algebra and $L^p$ norms (via \eqref{eq:norms});
\item freeness of $X_0,X_1,\dots$ (Voiculescu's theorem);
\item complex coefficients;
\item general $d$ (Proposition~\ref{prop:fuss}, Corollary~\ref{cor:alld});
\item the higher-chaos and unitary readings.
\end{itemize}

\paragraph{Python (\texttt{verify.py}, standard library).} Exact rational and Gaussian-rational
arithmetic. Moments are computed by two independent methods: sparse Fock-space operators, and
counting monochromatic non-crossing pairings. The two agree on all $9841$ words of length $\le8$
in three letters. The script checks:
\begin{itemize}
\item $\|W_d\|_4^4=d+1$ for $d=1..5$ by both methods and by the Fuss--Catalan formula;
\item all three normalisations for $d=1..5$;
\item the $2$-adic argument for $d<200$;
\item ratio $2$ on random real and complex first-chaos elements, including circular and
elliptic ones;
\item ratio $\ge2$ on random second-chaos elements, and ratio $2+(d-1)/n$ for sums of free copies;
\item the Haar-unitary formula.
\end{itemize}

\end{document}
